\section{Conclusions}
\label{sec:conclusions}

This work characterises scattered-light transients in LIGO O3 and O4a using
public data only, and asks how much of the physics survives that restriction.

\begin{enumerate}

\item \textbf{The morphology carries the mechanism.} Because the scattering
frequency follows $|v|$ rather than $v$, the arch repetition period measures
the driving motion directly, $\fdrive = 1/2\Tarch$, and the harmonic spacing
at the apex measures the peak surface velocity. Both are strain-only
measurements. Validated against injections, the arch period is recovered to
better than \SI{1}{\percent} and the harmonic spacing to within
\SI{7}{\percent}.

\item \textbf{Two public-data facts constrain what is possible, and neither is
documented elsewhere.} No Gravity Spy classification catalogue has been
released for O4a, and the O4a auxiliary release contains no seismometer,
accelerometer or suspension witness. Together they mean the O4a analysis must
be strain-only, and that the auxiliary correlation is necessarily an O3
result. Both were established by direct verification before any code was
written; the second is asserted in the test suite so that a future release
would be detected.

\item \textbf{The O3 auxiliary release supports more than expected.} Alongside
the complete ground-seismometer band-limited RMS family across the earthquake,
microseismic and anthropogenic bands, it contains
\code{L1:SUS-OM3\_M1\_DAMP\_L\_IN1\_DQ} at \SI{256}{\hertz} --- a genuine
optic-motion witness, and the only full-rate one in either release. This makes
a per-trigger test of $f_n = 2n|v|/\lambda$ possible from public data, which
was not obviously the case at the outset.

\item \textbf{The O3 auxiliary work is what licenses the O4a interpretation.}
Establishing in O3 that the trigger rate tracks ground motion, and that the
predictor explains individual arch trains, is what allows an O4a arch period
--- measured with no auxiliary data at all --- to be read as a statement about
the instrument's environment rather than as a description of a picture.

\item \textbf{Validation against known signals changed the analysis, twice.}
Capping the $Q$ range corrected a factor-of-twenty error in the measured arch
period of fast scattering, and raising the ridge threshold from $Z = 5.5$ to
$Z = 20$ removed a noise floor that was saturating the frequency-extent
features. Neither error was visible in the outputs; both were visible against
injections. The residual \SI{15}{\percent} high bias in maximum arch frequency
is reported rather than tuned away, and is used only where it cancels.

\end{enumerate}

The pipeline is public, resumable, and validated end to end. Its main
extension --- full O4a coverage --- is a matter of compute rather than method,
and the arrival of an O4a trigger catalogue would remove the largest remaining
systematic without changing anything else.

\subsection*{Acknowledgements}

This research has made use of data obtained from the Gravitational Wave Open
Science Center (\url{https://gwosc.org}), a service of the LIGO Laboratory, the
LIGO Scientific Collaboration, the Virgo Collaboration and KAGRA. LIGO
Laboratory and Advanced LIGO are funded by the United States National Science
Foundation, with contributions from the Science and Technology Facilities
Council of the United Kingdom, the Max-Planck-Society, and the State of
Niedersachsen/Germany.

This work uses the Gravity Spy machine-learning classification data release
\citep{GravitySpyData2021} and the GWOSC O3 and O4 auxiliary channel releases
\citep{GWOSCO3Aux,GWOSCO4Aux}. The Gravity Spy project is supported by the
United States National Science Foundation and relies on the contributions of
its citizen-science volunteers.

The analysis uses \code{gwpy} \citep{Macleod2021GWpy}, \code{numpy}
\citep{Harris2020NumPy}, \code{scipy} \citep{Virtanen2020SciPy},
\code{scikit-learn} \citep{Pedregosa2011} and \code{matplotlib}
\citep{Hunter2007}.

\subsection*{Data and code availability}

All input data are public: strain and auxiliary channels from GWOSC, and the
Gravity Spy classifications from Zenodo record 5649212. The analysis code,
configuration, verbatim channel lists and decision log are at
\url{https://github.com/Hardik-SSingh/ligo-scattered-light-o3-o4a}. Every derived
product is written with a provenance sidecar recording the configuration, the
random seed and the software versions used to produce it.
